\documentclass[11pt,a4paper]{article}
\usepackage[utf8]{inputenc}
\usepackage[T1]{fontenc}
\usepackage[french]{babel}
\usepackage{amsmath,amssymb,amsfonts,amsthm,mathrsfs,mathtools}
\usepackage[top=2cm,bottom=2cm,left=2cm,right=2cm]{geometry}
\usepackage[dvipsnames]{xcolor}
\usepackage{tcolorbox}
\tcbuselibrary{skins,breakable}
\usepackage{tikz}
\usetikzlibrary{calc,shapes,arrows.meta}
\usepackage{enumitem}
\usepackage{fancyhdr}
\usepackage{lastpage}
\usepackage{hyperref}

% ----------------------------------------------------
% Couleurs officielles du projet (Charte LLG / MPSI)
% ----------------------------------------------------
\definecolor{llgblue}{RGB}{20, 50, 115}
\definecolor{llglight}{RGB}{246, 249, 255}

\definecolor{emeraldframe}{RGB}{18, 110, 75}
\definecolor{emeraldback}{RGB}{242, 250, 246}
\definecolor{emeraldtitle}{RGB}{12, 80, 52}

\definecolor{prepaframe}{RGB}{30, 70, 135}
\definecolor{prepaback}{RGB}{246, 249, 255}
\definecolor{prepatitle}{RGB}{20, 50, 105}

\definecolor{alertframe}{RGB}{185, 75, 20}
\definecolor{alertback}{RGB}{255, 249, 245}
\definecolor{alerttitle}{RGB}{145, 55, 12}

% Configuration des hyperliens
\hypersetup{
    colorlinks=true,
    linkcolor=llgblue,
    urlcolor=llgblue,
    citecolor=llgblue
}

% En-têtes et pieds de page
\pagestyle{fancy}
\fancyhf{}
\renewcommand{\headrulewidth}{0.5pt}
\renewcommand{\footrulewidth}{0.4pt}
\fancyhead[L]{\small\textbf{Lycée Louis-le-Grand} -- 1\textsuperscript{ère} Spécialité Mathématiques (1E3)}
\fancyhead[C]{\small\textbf{DST : Dénombrement, Combinatoire \& Loi Binomiale}}
\fancyhead[R]{\small Juin 2022}
\fancyfoot[L]{\footnotesize LLG 2021/2022 -- Devoir Sur Table (1E3)}
\fancyfoot[C]{\footnotesize Énoncé \& Corrigé Intégral Rigoureux}
\fancyfoot[R]{\footnotesize Page \thepage\ / \pageref{LastPage}}

% ----------------------------------------------------
% Boîtes tcolorbox personnalisées
% ----------------------------------------------------
\newtcolorbox{correctionbox}[1][]{
  colback=emeraldback,
  colframe=emeraldframe,
  fonttitle=\bfseries\color{white},
  coltitle=white,
  colbacktitle=emeraldtitle,
  title={Corrigé officiel détaillé et barème indicatif #1},
  arc=2mm,
  boxrule=0.8pt,
  titlerule=0.5pt,
  breakable,
  before skip=0.4cm,
  after skip=0.5cm
}

\newtcolorbox{remarquebox}[1][]{
  colback=prepaback,
  colframe=prepaframe,
  fonttitle=\bfseries\small\color{white},
  coltitle=white,
  colbacktitle=prepatitle,
  title={Prolongement \& Regard Prépa / MPSI #1},
  arc=1.5mm,
  boxrule=0.6pt,
  titlerule=0.4pt,
  breakable,
  before skip=0.25cm,
  after skip=0.35cm
}

\newtcolorbox{alertebox}[1][]{
  colback=alertback,
  colframe=alertframe,
  fonttitle=\bfseries\small\color{white},
  coltitle=white,
  colbacktitle=alerttitle,
  title={Point de vigilance méthodologique \& Pièges fréquents #1},
  arc=1.5mm,
  boxrule=0.6pt,
  titlerule=0.4pt,
  breakable,
  before skip=0.25cm,
  after skip=0.35cm
}

% Commandes mathématiques
\newcommand{\N}{\mathbb{N}}
\newcommand{\Z}{\mathbb{Z}}
\newcommand{\Q}{\mathbb{Q}}
\newcommand{\R}{\mathbb{R}}
\newcommand{\C}{\mathbb{C}}
\newcommand{\Prob}{\mathbb{P}}
\newcommand{\E}{\mathbb{E}}
\newcommand{\V}{\mathbb{V}}
\newcommand{\un}{\mathbf{1}}

\begin{document}

% ====================================================
% EN-TÊTE DU DEVOIR
% ====================================================
\begin{center}
    {\LARGE \bfseries Lycée Louis-le-Grand -- Paris}\\[0.2cm]
    {\large Classe de Première 3 (1E3) -- Année scolaire 2021/2022}\\[0.3cm]
    \begin{tcolorbox}[colback=llglight,colframe=llgblue,arc=3mm,boxrule=1.2pt,center]
        \begin{center}
            {\Large \bfseries \color{llgblue} DEVOIR SUR TABLE DE MATHÉMATIQUES}\\[0.15cm]
            {\large \bfseries Dénombrement, Combinatoire Ensembliste, Identités Binomiales \& Tirages Sans Remise}\\[0.2cm]
            \small Date : Juin 2022 \quad $\vert$ \quad Durée : 2 heures \quad $\vert$ \quad Calculatrice autorisée selon la réglementation
        \end{center}
    \end{tcolorbox}
\end{center}

\vspace{0.3cm}

\begin{tcolorbox}[colback=white,colframe=llgblue!60,arc=1.5mm,boxrule=0.6pt]
\small\textbf{Consignes générales d'évaluation :}
\begin{itemize}[leftmargin=*,itemsep=1pt]
    \item La modélisation ensembliste (combinaisons, arrangements, p-uplets, permutations) doit être explicitée et justifiée pour chaque situation.
    \item Les démonstrations combinatoires par double comptage sont privilégiées par rapport aux calculs bruts de factorielles.
    \item Tout résultat numérique volumineux doit être donné sous forme exacte de coefficients binomiaux puis calculé.
\end{itemize}
\end{tcolorbox}

\vspace{0.4cm}

% ====================================================
% SUJET DU DEVOIR SURVEILLÉ (BARÈME /20)
% ====================================================
\begin{center}
    {\Large \bfseries \color{llgblue} \underline{Sujet du Devoir Surveillé (Barème sur 20 points)}}
\end{center}

\vspace{0.2cm}

% ----------------------------------------------------
% Exercice 1 (5 points)
% ----------------------------------------------------
\subsection*{Exercice 1 -- Mains de cartes \& Disjonctions d'intersection \hfill (5 points)}

On tire simultanément 5 cartes d'un jeu de 32 cartes (les valeurs vont du 7 à l'As dans les 4 enseignes : cœur, carreau, pique, trèfle).
Combien de mains différentes peut-on obtenir :
\begin{enumerate}
    \item Sans imposer de contraintes sur les cartes ?
    \item Contenant 5 carreaux ou 5 piques ?
    \item Contenant exactement 2 carreaux et 3 piques ?
    \item Contenant au moins un roi ?
    \item Contenant exactement 2 rois et 3 piques ? (\emph{Indication :} Dénombrer les mains sans le roi de pique et les mains avec le roi de pique, et montrer qu'il y a $1\,428$ mains possibles).
\end{enumerate}

\vspace{0.3cm}

% ----------------------------------------------------
% Exercice 2 (5 points)
% ----------------------------------------------------
\subsection*{Exercice 2 -- Dénombrement de parties \& Recouvrements \hfill (5 points)}

Soit $n \in \N^*$ et $E$ un ensemble fini possédant $n$ éléments. On désigne par $\mathscr{P}(E)$ l'ensemble des parties de $E$.
\begin{enumerate}
    \item Démontrer que $\operatorname{Card}(\mathscr{P}(E)) = 2^n$ par deux méthodes indépendantes (codage binaire et formule du binôme).
    \item Déterminer le nombre de couples $(A, B) \in (\mathscr{P}(E))^2$ tels que $B \subset A$.
    \item Déterminer le nombre de couples $(A, B) \in (\mathscr{P}(E))^2$ tels que $A \cap B = \emptyset$.
    \item Déterminer le nombre de couples $(A, B) \in (\mathscr{P}(E))^2$ tels que $A \cup B = E$.
    \item Déterminer le nombre de triplets $(A, B, C) \in (\mathscr{P}(E))^3$ tels que $A \cup B \cup C = E$.
\end{enumerate}

\vspace{0.3cm}

% ----------------------------------------------------
% Exercice 3 (5 points)
% ----------------------------------------------------
\subsection*{Exercice 3 -- Formule du pion, Identité de Vandermonde \& Inversion de Pascal \hfill (5 points)}

Soient $i, k, n$ des entiers naturels tels que $i \le k \le n$. Par convention, $0^0 = 1$.
\begin{enumerate}
    \item Démontrer par un dénombrement combinatoire (en considérant le choix d'un comité puis d'un sous-comité) la formule du pion :
    \[
    \binom{n}{k} \binom{k}{i} = \binom{n}{i} \binom{n - i}{k - i}
    \]
    \item Soit $q \in \N$ avec $q \le p \le n$. Démontrer par un dénombrement que :
    \[
    \binom{n}{p} = \sum_{k=0}^q \binom{q}{k} \binom{n - q}{p - k}
    \]
    En déduire la somme remarquable : $\displaystyle\sum_{k=0}^n \binom{n}{k}^2 = \binom{2n}{n}$.
    \item Montrer que :
    \[
    \sum_{k=i}^n (-1)^{n - k} \binom{n}{k} \binom{k}{i} = \binom{n}{i} 0^{n - i} = \begin{cases} 1 & \text{si } i = n \\ 0 & \text{si } i < n \end{cases}
    \]
    \item Soient $(u_n)_{n \in \N}$ et $(v_n)_{n \in \N}$ deux suites réelles liées par :
    \[
    \forall n \in \N, \quad u_n = \sum_{k=0}^n \binom{n}{k} v_k
    \]
    Démontrer la formule d'\emph{inversion de Pascal} :
    \[
    \forall n \in \N, \quad v_n = \sum_{k=0}^n (-1)^{n - k} \binom{n}{k} u_k
    \]
\end{enumerate}

\vspace{0.3cm}

% ----------------------------------------------------
% Exercice 4 (5 points)
% ----------------------------------------------------
\subsection*{Exercice 4 -- Temps d'attente sans remise dans une urne bicolore \hfill (5 points)}

Soit $n \in \N^*$. Une urne contient $n$ boules blanches et $2$ boules noires indiscernables au toucher (soit $n+2$ boules au total).
On effectue des tirages successifs sans remise de toutes les boules de l'urne.
On note $X$ la variable aléatoire égale au rang d'apparition de la première boule blanche.

\begin{enumerate}
    \item Justifier que les valeurs possibles prises par $X$ sont $X(\Omega) = \{1, 2, 3\}$.
    \item Déterminer la loi de probabilité de la variable aléatoire $X$.
    \item Démontrer que l'espérance mathématique de $X$ est donnée par :
    \[
    \E(X) = \frac{n + 3}{n + 1} = 1 + \frac{2}{n + 1}
    \]
    \item Étudier la limite de $\E(X)$ quand $n \to +\infty$ et interpréter ce résultat.
\end{enumerate}

\newpage

% ====================================================
% CORRIGÉ OFFICIEL DU DEVOIR SURVEILLÉ
% ====================================================
\begin{center}
    {\Large \bfseries \color{emeraldtitle} \underline{Corrigé Intégral Officiel \& Approfondissements Prépa}}
\end{center}

\vspace{0.3cm}

% ----------------------------------------------------
% Correction Exercice 1
% ----------------------------------------------------
\subsection*{Correction détaillée de l'Exercice 1 \hfill (5 points)}

\begin{correctionbox}[-- Analyse combinatoire des mains de cartes]
Une main de cartes est un sous-ensemble (sans ordre et sans répétition) de 5 cartes choisies parmi 32.
\begin{enumerate}
    \item \textbf{Nombre de mains sans contraintes :}
    Il s'agit du nombre de combinaisons de 5 cartes parmi 32 :
    \[
    \operatorname{Card}(\Omega) = \binom{32}{5} = \frac{32 \times 31 \times 30 \times 29 \times 28}{5 \times 4 \times 3 \times 2 \times 1} = \mathbf{201\,376\text{ mains}}
    \]

    \item \textbf{Mains contenant 5 carreaux ou 5 piques :}
    Les événements \og 5 carreaux \fg{} et \og 5 piques \fg{} sont incompatibles.
    Chaque couleur compte 8 cartes. Le nombre de mains est donc :
    \[
    \binom{8}{5} + \binom{8}{5} = 2 \times \frac{8 \times 7 \times 6}{3 \times 2 \times 1} = 2 \times 56 = \mathbf{112\text{ mains}}
    \]

    \item \textbf{Mains contenant 2 carreaux et 3 piques :}
    On choisit indépendamment 2 carreaux parmi les 8 carreaux et 3 piques parmi les 8 piques :
    \[
    \binom{8}{2} \times \binom{8}{3} = \frac{8 \times 7}{2} \times 56 = 28 \times 56 = \mathbf{1\,568\text{ mains}}
    \]

    \item \textbf{Mains contenant au moins un roi :}
    Considérons l'événement contraire : \og la main ne contient aucun roi \fg{}.
    Il y a 4 rois et $32 - 4 = 28$ cartes non-rois.
    Le nombre de mains sans aucun roi est $\binom{28}{5} = \frac{28 \times 27 \times 26 \times 25 \times 24}{120} = 98\,280$.
    Par soustraction :
    \[
    N(\text{au moins un roi}) = \binom{32}{5} - \binom{28}{5} = 201\,376 - 98\,280 = \mathbf{103\,096\text{ mains}}
    \]

    \item \textbf{Mains contenant 2 rois et 3 piques :}
    Le roi de pique ($R_\spadesuit$) jouant un double rôle (à la fois roi et pique), partitionnons selon son appartenance à la main :
    \begin{itemize}[label=$\bullet$]
        \item \emph{Cas 1 : La main ne contient PAS le roi de pique ($R_\spadesuit \notin \text{main}$) :}
        La main doit contenir 2 rois parmi les 3 autres rois ($\binom{3}{2} = 3$ choix) et 3 piques parmi les 7 autres piques ($\binom{7}{3} = 35$ choix). La main est alors complète (5 cartes) :
        \[
        N_1 = \binom{3}{2} \times \binom{7}{3} = 3 \times 35 = \mathbf{105\text{ mains}}
        \]
        \item \emph{Cas 2 : La main CONTIENT le roi de pique ($R_\spadesuit \in \text{main}$) :}
        $R_\spadesuit$ fournit déjà 1 roi et 1 pique. Il faut donc adjoindre à la main :
        \begin{itemize}
            \item $1$ autre roi parmi les 3 rois restants non-piques ($\binom{3}{1} = 3$ choix) ;
            \item $2$ autres piques parmi les 7 piques restants non-rois ($\binom{7}{2} = 21$ choix) ;
            \item $1$ carte restante qui n'est NI un roi NI un pique. Le nombre de cartes disponibles pour ce choix est $32 - 4 \text{ (rois)} - 7 \text{ (piques restants)} = 21$ cartes ordinaires ($\binom{21}{1} = 21$ choix).
        \end{itemize}
        \[
        N_2 = \binom{3}{1} \times \binom{7}{2} \times \binom{21}{1} = 3 \times 21 \times 21 = \mathbf{1\,323\text{ mains}}
        \]
    \end{itemize}
    Par disjonction des cas, le nombre total de mains est :
    \[
    N_{\text{total}} = N_1 + N_2 = 105 + 1\,323 = \mathbf{1\,428\text{ mains}}
    \]
\end{enumerate}
\end{correctionbox}

\vspace{0.4cm}

% ----------------------------------------------------
% Correction Exercice 2
% ----------------------------------------------------
\subsection*{Correction détaillée de l'Exercice 2 \hfill (5 points)}

\begin{correctionbox}[-- Dénombrement de parties \& Recouvrements]
Soit $E$ un ensemble de cardinal $n$.
\begin{enumerate}
    \item \textbf{Démonstrations de $\operatorname{Card}(\mathscr{P}(E)) = 2^n$ :}
    \begin{itemize}[label=$\bullet$]
        \item \emph{Méthode 1 (Fonctions caractéristiques) :} Une partie $A \subset E$ est entièrement caractérisée par son indicatrice $\un_A : E \to \{0, 1\}$. Il y a donc une bijection entre $\mathscr{P}(E)$ et $\{0, 1\}^E$.
        Or $\operatorname{Card}(\{0, 1\}^E) = 2^{\operatorname{Card}(E)} = 2^n$.
        \item \emph{Méthode 2 (Somme des combinaisons \& Binôme) :} $\mathscr{P}(E) = \bigcup_{k=0}^n \{ A \subset E \mid \operatorname{Card}(A) = k \}$.
        Puisque cette union est disjointe :
        \[
        \operatorname{Card}(\mathscr{P}(E)) = \sum_{k=0}^n \binom{n}{k} = \sum_{k=0}^n \binom{n}{k} 1^k 1^{n-k} = (1 + 1)^n = 2^n
        \]
    \end{itemize}

    \item \textbf{Couples $(A, B)$ tels que $B \subset A$ :}
    Pour chaque élément $x \in E$, trois situations mutuellement exclusives sont possibles :
    $x \in B$ (ce qui implique $x \in A$) ; $x \in A \setminus B$ ; $x \notin A$ (donc $x \notin B$).
    Il y a donc 3 choix indépendants pour chacun des $n$ éléments de $E$.
    Le nombre de couples est donc :
    \[
    \mathbf{3^n}
    \]
    \emph{Vérification par sommation :} $\sum_{k=0}^n \binom{n}{k} 2^k = (1 + 2)^n = 3^n$.

    \item \textbf{Couples $(A, B)$ disjoints ($A \cap B = \emptyset$) :}
    Pour chaque élément $x \in E$, il y a 3 choix : $x \in A$ (donc $x \notin B$) ; $x \in B$ (donc $x \notin A$) ; $x \notin A \cup B$.
    Le nombre de couples est de même :
    \[
    \mathbf{3^n}
    \]

    \item \textbf{Couples $(A, B)$ recouvrant $E$ ($A \cup B = E$) :}
    Par passage aux complémentaires, $A \cup B = E \iff \overline{A} \cap \overline{B} = \emptyset$.
    L'application $(A, B) \mapsto (\overline{A}, \overline{B})$ étant une bijection de $(\mathscr{P}(E))^2$ dans lui-même, il y a autant de couples recouvrant $E$ que de couples disjoints :
    \[
    \mathbf{3^n}
    \]

    \item \textbf{Triplets $(A, B, C)$ tels que $A \cup B \cup C = E$ :}
    Pour que l'union soit $E$, aucun élément de $E$ ne doit appartenir au complémentaire commun $\overline{A} \cap \overline{B} \cap \overline{C}$.
    Pour chaque élément $x \in E$, le sous-ensemble de $\{A, B, C\}$ contenant $x$ est un sous-ensemble \textbf{non vide} de $\{A, B, C\}$.
    Il y a $2^3 - 1 = 8 - 1 = 7$ choix possibles d'appartenance pour chaque élément $x \in E$.
    Puisque les choix sont indépendants pour les $n$ éléments :
    \[
    \mathbf{7^n\text{ triplets}}
    \]
\end{enumerate}
\end{correctionbox}

\vspace{0.4cm}

% ----------------------------------------------------
% Correction Exercice 3
% ----------------------------------------------------
\subsection*{Correction détaillée de l'Exercice 3 \hfill (5 points)}

\begin{correctionbox}[-- Formule du pion, Vandermonde \& Inversion de Pascal]
\begin{enumerate}
    \item \textbf{Formule du pion par double comptage :}
    Considérons la formation d'un comité $A$ de $k$ personnes parmi $n$, puis d'un sous-comité directeur $B$ de $i$ personnes au sein de $A$ ($i \le k \le n$).
    \begin{itemize}[label=$\bullet$]
        \item Méthode 1 : On choisit d'abord les $k$ membres du comité parmi $n$ ($\binom{n}{k}$ choix), puis les $i$ dirigeants parmi les $k$ ($\binom{k}{i}$ choix) : $\binom{n}{k} \binom{k}{i}$ manières.
        \item Méthode 2 : On choisit directement les $i$ dirigeants parmi $n$ ($\binom{n}{i}$ choix), puis on complète le comité en choisissant les $k - i$ membres restants parmi les $n - i$ personnes disponibles ($\binom{n-i}{k-i}$ choix) : $\binom{n}{i} \binom{n-i}{k-i}$ manières.
    \end{itemize}
    Les deux dénombrements comptabilisant les mêmes configurations $(B \subset A)$ :
    \[
    \mathbf{\binom{n}{k} \binom{k}{i} = \binom{n}{i} \binom{n - i}{k - i}}
    \]

    \item \textbf{Identité de Vandermonde :}
    Soit un ensemble de $n$ personnes composé de $q$ femmes et de $n - q$ hommes.
    On souhaite former une délégation de $p$ personnes. Le nombre total de délégations est $\binom{n}{p}$.
    En partitionnant selon le nombre $k$ de femmes dans la délégation ($0 \le k \le q$), on choisit $k$ femmes parmi $q$ ($\binom{q}{k}$) et les $p - k$ hommes restants parmi $n - q$ ($\binom{n-q}{p-k}$).
    Par additivité sur les cas disjoints :
    \[
    \mathbf{\binom{n}{p} = \sum_{k=0}^q \binom{q}{k} \binom{n - q}{p - k}}
    \]
    Pour $p = q = n$ dans un ensemble de $2n$ éléments divisé en deux groupes de taille $n$ :
    \[
    \binom{2n}{n} = \sum_{k=0}^n \binom{n}{k} \binom{n}{n - k} = \sum_{k=0}^n \binom{n}{k} \binom{n}{k} = \mathbf{\sum_{k=0}^n \binom{n}{k}^2}
    \]

    \item \textbf{Calcul de la somme alternée :}
    En substituant la formule du pion établie en 1 :
    \[
    S = \sum_{k=i}^n (-1)^{n - k} \binom{n}{k} \binom{k}{i} = \sum_{k=i}^n (-1)^{n - k} \binom{n}{i} \binom{n - i}{k - i}
    \]
    Le facteur $\binom{n}{i}$ ne dépend pas de $k$, on peut le factoriser :
    \[
    S = \binom{n}{i} \sum_{k=i}^n (-1)^{n - k} \binom{n - i}{k - i}
    \]
    Effectuons le changement d'indice $m = k - i$ (alors $n - k = (n - i) - m$) :
    \[
    S = \binom{n}{i} \sum_{m=0}^{n - i} (-1)^{(n - i) - m} \binom{n - i}{m}
    \]
    Reconnaissons le développement du binôme de Newton pour $(1 - 1)^{n - i}$ :
    \[
    S = \binom{n}{i} (1 - 1)^{n - i} = \mathbf{\binom{n}{i} 0^{n - i} = \begin{cases} 1 & \text{si } i = n \\ 0 & \text{si } i < n \end{cases}}
    \]

    \item \textbf{Formule d'inversion de Pascal :}
    Injectons $u_k = \sum_{j=0}^k \binom{k}{j} v_j$ dans la somme proposée :
    \[
    \sum_{k=0}^n (-1)^{n - k} \binom{n}{k} u_k = \sum_{k=0}^n (-1)^{n - k} \binom{n}{k} \left( \sum_{j=0}^k \binom{k}{j} v_j \right)
    \]
    Intervertissons les deux sommations (domaine $0 \le j \le k \le n$) :
    \[
    \sum_{j=0}^n v_j \left( \sum_{k=j}^n (-1)^{n - k} \binom{n}{k} \binom{k}{j} \right)
    \]
    D'après la question 3, la somme intérieure vaut $\delta_{nj}$ (elle s'annule pour tout $j < n$ et vaut 1 pour $j = n$).
    Il ne reste donc que le terme $j = n$ :
    \[
    v_n \times 1 = \mathbf{v_n}
    \]
\end{enumerate}
\end{correctionbox}

\vspace{0.4cm}

% ----------------------------------------------------
% Correction Exercice 4
% ----------------------------------------------------
\subsection*{Correction détaillée de l'Exercice 4 \hfill (5 points)}

\begin{correctionbox}[-- Rang d'apparition du premier succès sans remise]
\textbf{1. Univers image $X(\Omega) = \{1, 2, 3\}$ :}
L'urne contient $n$ boules blanches ($B$) et exactement $2$ boules noires ($N$).
Le pire des scénarios pour tirer une boule blanche est d'extraire successivement les deux seules boules noires aux deux premiers tirages ($N_1 \cap N_2$). La troisième boule extraite sera alors obligatoirement blanche ($B_3$).
Par conséquent, la première boule blanche apparaît nécessairement au premier, au deuxième ou au troisième tirage :
\[
\mathbf{X(\Omega) = \{1, 2, 3\}}
\]

\textbf{2. Loi de probabilité de $X$ :}
\begin{itemize}[label=$\bullet$]
    \item \textbf{$X = 1$ :} La première boule tirée est blanche :
    \[
    P(X = 1) = \mathbf{\frac{n}{n + 2}}
    \]
    \item \textbf{$X = 2$ :} La première est noire et la deuxième est blanche :
    \[
    P(X = 2) = P(N_1) \times P_{N_1}(B_2) = \frac{2}{n + 2} \times \frac{n}{n + 1} = \mathbf{\frac{2n}{(n + 1)(n + 2)}}
    \]
    \item \textbf{$X = 3$ :} Les deux premières sont noires et la troisième est blanche :
    \[
    P(X = 3) = \frac{2}{n + 2} \times \frac{1}{n + 1} \times \frac{n}{n} = \mathbf{\frac{2}{(n + 1)(n + 2)}}
    \]
\end{itemize}
\emph{Vérification de la somme :}
\[
\frac{n(n + 1) + 2n + 2}{(n + 1)(n + 2)} = \frac{n^2 + 3n + 2}{(n + 1)(n + 2)} = \frac{(n + 1)(n + 2)}{(n + 1)(n + 2)} = 1
\]

\textbf{3. Calcul de l'espérance mathématique $\E(X)$ :}
Par définition de l'espérance :
\[
\E(X) = 1 \times P(X = 1) + 2 \times P(X = 2) + 3 \times P(X = 3)
\]
\[
\E(X) = 1 \times \frac{n}{n + 2} + 2 \times \frac{2n}{(n + 1)(n + 2)} + 3 \times \frac{2}{(n + 1)(n + 2)}
\]
Réduisons au même dénominateur $(n + 1)(n + 2)$ :
\[
\E(X) = \frac{n(n + 1) + 4n + 6}{(n + 1)(n + 2)} = \frac{n^2 + n + 4n + 6}{(n + 1)(n + 2)} = \frac{n^2 + 5n + 6}{(n + 1)(n + 2)}
\]
Factorisons le numérateur $n^2 + 5n + 6 = (n + 2)(n + 3)$ :
\[
\E(X) = \frac{(n + 2)(n + 3)}{(n + 1)(n + 2)} = \mathbf{\frac{n + 3}{n + 1}} = \mathbf{1 + \frac{2}{n + 1}}
\]

\textbf{4. Limite et interprétation probabiliste :}
Quand le nombre de boules blanches $n$ devient très grand ($n \to +\infty$) :
\[
\lim_{n \to +\infty} \E(X) = \lim_{n \to +\infty} \left( 1 + \frac{2}{n + 1} \right) = \mathbf{1}
\]
\textbf{Interprétation :} Lorsque l'urne est massivement dominée par les boules blanches, la probabilité d'obtenir une boule blanche dès le tout premier tirage tend vers 1 ($P(X=1) = \frac{n}{n+2} \to 1$). Le temps d'attente moyen pour observer la première blanche converge donc naturellement vers 1 tirage.
\end{correctionbox}

\end{document}
